where @@M0@@ is a
smooth morphism. One may also suppose @@M1@@, @@M2@@, @@M3@@ of constant rank, as well as the conormal sheaves of @@M4@@ and @@M5@@.

From the commutative diagram (in which all morphisms are smooth):
@@M6@@
one deduces:
\sgaStableTarget{sga6.e08.eq.1-8-1}@@M7@@

Let @@M8@@; then the rank of @@M9@@ is equal to that of @@M10@@. Furthermore, @@M11@@ defines a section @@M12@@, which is a regular immersion of codimension equal to the rank of @@M13@@ (EGA IV 17.10.4). If @@M14@@ is the immersion obtained from @@M15@@ by the base change @@M16@@, we have @@M17@@ and consequently (VII 1.7):
\sgaStableTarget{sga6.e08.eq.1-8-2}@@M18@@

Comparing (1.8.1) and (1.8.2) gives the proposition.

\sgaStableTarget{sga6.e08.def.1-9}\textbf{Definition 1.9.} \uline{The integer} @@M19@@ \uline{of Proposition} 1.8 \uline{is called the} virtual relative dimension \uline{of} @@M20@@ \uline{over} @@M21@@ (\uline{or of} @@M22@@) \uline{at the point} @@M23@@.

The virtual relative dimension of @@M24@@ over @@M25@@ is therefore a locally constant function on @@M26@@, with integer values (positive or negative). If @@M27@@ is a smooth morphism of relative dimension @@M28@@, its virtual relative dimension is @@M29@@. If @@M30@@ is a regular immersion of codimension @@M31@@, its virtual relative dimension is @@M32@@.

\sgaStableTarget{sga6.e08.prop.1-10}\textbf{Proposition 1.10.} \uline{Let} @@M33@@ \uline{and} @@M34@@ \uline{be two complete-intersection morphisms}, \uline{and let} @@M35@@ \uline{be a point of} @@M36@@. \uline{Then}:
@@M37@@
